\documentclass[11pt,a4paper]{article}

% =======================
% Packages
% =======================
\usepackage[utf8]{inputenc}
\usepackage[T1]{fontenc}
\usepackage[french]{babel}
\usepackage{lmodern}
\usepackage{microtype}
\usepackage{geometry}
\usepackage{graphicx}
\usepackage{booktabs}
\usepackage{hyperref}
\usepackage{float}
\usepackage{enumitem}
\usepackage{titlesec}
\usepackage{amsmath}
\usepackage{amssymb}
\usepackage{xcolor}

% =======================
% Mise en page compacte (Note de synthèse 4 pages max)
% =======================
\geometry{margin=1.4cm}
\setlength{\parindent}{0pt}
\setlength{\parskip}{4pt}
\setlist[itemize]{leftmargin=*, nosep}
\setlist[enumerate]{leftmargin=*, nosep}

\titlespacing*{\section}{0pt}{0.6em}{0.2em}
\titlespacing*{\subsection}{0pt}{0.35em}{0.12em}

\setlength{\textfloatsep}{7pt}
\setlength{\floatsep}{6pt}
\setlength{\intextsep}{6pt}

\hypersetup{colorlinks=true, linkcolor=blue, urlcolor=blue}

% =======================
% Titre
% =======================
\title{\vspace{-1.2cm}\textbf{Note de Synthèse -- Classification Hiérarchique de Catalogue}\\
\large Architecture globale, modèles supervisés et enrichissement LLM local}
\author{\small Salah Eddine NEJJARI \;|\; Tao DANG TRAN \;|\; Thomas FAVRE \;|\; Yohan ANDRIAMAMPIONONA}
\date{\small Projet Stat'App -- ENSAE Paris \;|\; Mise à jour : Mai 2026}

\begin{document}
\maketitle
\vspace{-0.9cm}

\small
\textbf{Résumé.} Ce projet vise à projeter un catalogue massif de 128\,000 produits vers la taxonomie complexe de Google (5440 nœuds, profondeur max 7). Initialement exploré via une approche algorithmique non-supervisée (Min-Cost Max-Flow), le projet a basculé vers un \textbf{Pipeline Hiérarchique à Double Modèle (Dual Embedding Pipeline)}. Le Niveau 1 (bénéficiant de Ground Truth) utilise un apprentissage de représentation poussé suivi d'un XGBoost. Les Niveaux 2 et inférieurs s'appuient sur un enrichissement par \textbf{Large Language Model exécuté localement (Apple MLX)}, un système multi-prototypes, et une inférence hiérarchique bornée par un Search Beam.

\vspace{0.2cm}
\textbf{Contexte Industriel (Critéo) :} Ce projet est un cas d'étude pédagogique exploratoire proposé par Critéo, sans objectif de mise en production. Acteur majeur de la publicité en ligne, Critéo opère le placement de bannières numériques personnalisées en ingérant des catalogues massifs en provenance de nombreux annonceurs. Classer ces produits finement et avec un taux d'erreur minime est critique pour l'entreprise : cette compréhension sémantique de l'inventaire publicitaire permet aux algorithmes d'ad-tech de déterminer exactement quelle annonce afficher à quel consommateur pour maximiser la pertinence et les conversions.
\vspace{0.2cm}

% ==============================================================================
\section{1. Défis et Approche Initiale : Apprentissage sous contrainte (Phase 1)}
% ==============================================================================
Face à la volumétrie des données, notre approche initiale cherchait à limiter la dépendance à l'annotation manuelle. Nous avons modélisé l'assignation au Niveau 1 (34 catégories) grâce à un réseau de transport biparti. Chaque produit est converti en vecteur dense via \texttt{all-MiniLM-L6-v2} ; chaque catégorie devient un barycentre vectoriel de sa branche cible associée. 

L'algorithme \textit{Min-Cost Max-Flow} (résolu par OR-Tools) assignait globalement chaque texte en forçant le respect strict des quotas d'occurrence par racine. Cette phase a démontré qu'une régularisation macroscopique stabilisait les prédictions en Zero-shot, mais restait un palliatif aux frontières sémantiques très floues des réels usages d'e-commerce.

% ==============================================================================
\section{2. Apprentissage Supervisé au Level 1 : Triplet Loss et XGBoost}
% ==============================================================================
L'acquisition d'un set de \textit{Ground Truth} complet sur le Niveau 1 a permis une refonte du système selon une logique séquentielle supervisée :
\begin{enumerate}
    \item \textbf{Création de l'Espace Latent (Batch Hard Triplet Loss) :} Le modèle sémantique de base a été fine-tuné via une fonction  $\mathcal{L}_{\text{BHTL}}$ pour rapprocher les produits virtuellement cibles d'une même catégorie et fortement rejeter les négatifs les plus proches. 
    \item \textbf{Gestion GPU locale (MPS) :} Les limitations structurelles à la mémoire unifiée des puces Apple nous ont contraints à une forte ingénierie : micro-batching ($B=4, 8$), limitation de tokens à 128 et sous-échantillonnage maîtrisé (30k lignes).
    \item \textbf{Frontières non-linéaires (XGBoost) :} Un classifieur eXtreme Gradient Boosting capte les subtiles non-linéarités de ces embeddings en 384 dimensions, portant l'Accuracy globale à \textbf{91\%} (+26 points par rapport au système par Quotas, avec des bonds spectaculaires sur les \textit{Arts \& Entertainment}, de 13\% à 68\%).
\end{enumerate}

% ==============================================================================
\section{3. Level 2+ : Enrichissement LLM et Stratégie Multi-Prototypes}
% ==============================================================================
À partir du Level 2, le modèle fait face à l'impossibilité d'utiliser l'apprentissage supervisé, faute de vérité terrain, alors qu'il y a plus de 5000 catégories possibles.
Nous avons opté pour une \textbf{stratégie ascendante de Zero-shot augmentée}, pilotée par des modèles génératifs textuels et de la recherche géométrique.

\subsection{Génération synthétique sur environnement contraint (MLX)}
Afin de donner de l'expressivité sémantique aux simples labels Google, un modèle linguistique reconstruit massivement le profil de chaque sous-catégorie. Une des majeures difficultés de 2026 fut technique : l'utilisation du framework torch \textit{MPS} local (Mac) créait d'importants dépassements mémoire ("OOM"). 

L'ingénierie s'est alors tournée vers \textbf{la stack MLX (Apple Silicon native)}. Nous traitons les descriptions de manière dynamique via un modèle 4-bit, \textbf{Qwen2.5-7B-Instruct-4bit}. Le backend MLX, couplé à des requêtes \textit{batchées} en JSON restreint (ne renvoyant que des variantes lexicales pures), permet de recréer massivement des textes de références synthétiques sans déporter la donnée en inférence sur des API tierces.

\subsection{Local Scoring, Graphe NetworkX et Multi-Prototypes}
Le L1 figé par le XGBoost, un algorithme restrictif (\texttt{NetworkX}) dresse un sous-graphe des enfants possibles. 
Au lieu de comparer un produit à un centre de gravité fixe, le produit interroge une liste de \textbf{différents prototypes textuels cibles par catégorie} (nom de catégorie, sous-descendants possibles générés par le LLM, et variantes lexicales). Une aggrégation (ex: \texttt{mean\_top\_k}) calcule un score de parenté à partir d'un Embedding expert en \textit{Semantic Textual Similarity} (BGE-Small).

% ==============================================================================
\section{4. Fiabilisation du parcours via Beam Search et MLOps}
% ==============================================================================

\textbf{Soft Hierarchical Decoding (Beam Search)} \\
Un enjeu majeur du parcours hiérarchique est la cascade d'erreur : se tromper au Level 1 condamne entièrement les ramifications inférieures (Hard decoding). Pour contourner cela, notre implémentation supporte maintenant un \textbf{Beam Search borné}. En cas d'ambiguïté forte sur les n\oe{}uds parents, le système conserve les $B$ meilleures ramifications. Une fonction de combinaison hybride, $\alpha\cdot\mathbb{P}(\text{L1}) + (1-\alpha)\cdot S_{\text{cos}}(\text{L2}_n)$, permet ainsi aux traits éminemment évidents du plus bas niveau de taxonomie de sauver une erreur d'alignement à la racine.

\textbf{Architecturation Industrielle}\\
L'ensemble de ces briques est aujourd'hui découplé des vieux notebooks de recherche. Le code se répartit en un package \texttt{src/} stable et des \texttt{scripts/} permettant une exécution sérielle (preprocessing, scripts de fine-tuning Kaggle automatisés, indexation multi-prototypes). Une interface d'annotation interactive par \textit{Terminal} a été créée dans le seul but de réaliser de l'Active Learning humain accéléré, afin de construire in fine un outil d'évaluation rigoureux (Evaluation Set) pour les couches profondes (L2+).

% ==============================================================================
\section{Conclusion et Prochaines Étapes}
% ==============================================================================
Le pipeline a prouvé que la scission du processus -- \textbf{1. Apprentissage lourd et verrouillage structurel de la racine} / \textbf{2. Exploration Zéro-Shot sur graphes boostés par prototypes LLM MLX} -- est la seule approche soutenable pour un catalogue d'e-commerce massif évoluant rapidement.
Les prochaines instances du projet se concentreront sur la perturbation asymétrique de l'espace métrique de recherche euclidien ($L_2$) vis-à-vis du classique cosinus de similarité, et sur le bouclage des dernières métriques de robustesse sur l'échantillon réel L2 annoté.

\end{document}